\ExamPart{Trắc nghiệm trả lời ngắn}{2}{Thí sinh trả lời từ câu 1 đến câu 4.}

\NQuestion Một cửa hàng bán một loại sản phẩm với giá $250$ nghìn đồng mỗi chiếc. Chi phí sản xuất $x$ chiếc là $C(x)=x^2+50x+600 \quad (\text{nghìn đồng})$. Hỏi cửa hàng phải bán ít nhất bao nhiêu chiếc để không bị lỗ?
\AnswerLine{$4$}
\SolutionBlock{Doanh thu là $R(x)=250x$ (nghìn đồng). Điều kiện không lỗ là
\[
R(x)\ge C(x) \Leftrightarrow 250x\ge x^2+50x+600.
\]
Suy ra
\[
x^2-200x+600\le0.
\]
Giải phương trình $x^2-200x+600=0$ được hai nghiệm
\[
x=100\pm10\sqrt{94}.
\]
Vì $100-10\sqrt{94}\approx3{,}05$ và $x$ là số nguyên dương, nên số chiếc ít nhất phải bán là $4$.}

\NQuestion Trong mặt phẳng tọa độ $Oxy$, một điểm giao hàng ở vị trí $M(3;4)$. Kho hàng nằm trên đường thẳng $d:x+2y-8=0$. Tính quãng đường ngắn nhất từ $M$ đến kho hàng.
\AnswerLine{$\dfrac{3\sqrt5}{5}$}
\SolutionBlock{Khoảng cách từ $M(3;4)$ đến đường thẳng $d$ là
\[
d(M,d)=\dfrac{|3+2\cdot4-8|}{\sqrt{1^2+2^2}}=\dfrac{3}{\sqrt5}=\dfrac{3\sqrt5}{5}.
\]
Do đó quãng đường ngắn nhất là $\dfrac{3\sqrt5}{5}$.}

\NQuestion Phương trình $x^4-5x^2+4=0$ có bao nhiêu nghiệm thực phân biệt?
\AnswerLine{$4$}
\SolutionBlock{Đặt $t=x^2\ (t\ge0)$, ta được
\[
t^2-5t+4=0 \Leftrightarrow (t-1)(t-4)=0.
\]
Suy ra $t=1$ hoặc $t=4$, do đó $x=\pm1,\pm2$. Vậy có $4$ nghiệm thực phân biệt.}

\NQuestion Có bao nhiêu số nguyên $m$ để phương trình $|x^2-4x+3|=m$ có đúng $4$ nghiệm thực phân biệt?
\AnswerLine{$0$}
\SolutionBlock{Ta có
\[
x^2-4x+3=(x-2)^2-1.
\]
Phương trình đã cho tương đương với
\[
(x-2)^2-1=m \quad \text{hoặc} \quad (x-2)^2-1=-m,
\]
hay
\[
(x-2)^2=m+1 \quad \text{hoặc} \quad (x-2)^2=1-m.
\]
Vì $|x^2-4x+3|\ge0$ nên $m\ge0$.
\begin{itemize}[leftmargin=1.5em]
\item Nếu $m=0$ thì phương trình trở thành $x^2-4x+3=0$, có $2$ nghiệm phân biệt.
\item Nếu $m=1$ thì ta được $(x-2)^2=2$ hoặc $(x-2)^2=0$, có $3$ nghiệm phân biệt.
\item Nếu $m>1$ thì phương trình $(x-2)^2=1-m$ vô nghiệm, còn $(x-2)^2=m+1$ có đúng $2$ nghiệm phân biệt.
\end{itemize}
Vậy không có số nguyên $m$ nào để phương trình có đúng $4$ nghiệm thực phân biệt.}
